\section{Record codec and cut accounting}
\label{supp:codec}

This section gives the field-level definitions behind the charged model of
the main text.  Every lower bound there holds for any self-delimiting
serialization whose control field records the emitted-prefix length; the
definitions below fix the serialization used by the upper bounds and by
the instrumented implementation.

\begin{definition}[Sparse formal-parameter serialization]
\label{def:parmodule}
Fix formal angle symbols $\theta_1,\ldots,\theta_M$.  A parameter is stored
in normalized sparse form
\begin{equation*}
\begin{gathered}
 \lambda=\displaystyle\sum_{\ell=1}^{s}c_\ell\theta_{j_\ell}
 +\dfrac{a}{d}\pi,\\
 j_1<\cdots<j_s,\quad c_\ell,a\in\mathbb Z,\quad d\geq1,
\end{gathered}
\end{equation*}
with zero coefficients omitted and $\gcd(a,d)=1$.  Let
$\mathsf{u}(z)$ and $\mathsf{s}(z)$ be fixed self-delimiting codes for
nonnegative and signed integers.  The executable parameter serialization is
\begin{equation}
\label{eq:precision-charge}
\begin{aligned}
 \mathsf{ser}_{\mathrm{par}}(\lambda)
 &=\mathsf{u}(s)\prod_{\ell=1}^{s}
   \bigl(\mathsf{u}(j_\ell)\mathsf{s}(c_\ell)\bigr)\\
 &\quad\cdot\mathsf{s}(a)\mathsf{u}(d),\\
 \pcharge(\lambda)&:=|\mathsf{ser}_{\mathrm{par}}(\lambda)|.
\end{aligned}
\end{equation}
Thus one nonzero formal coefficient costs $O(\log M)$ rather than
$\Omega(M)$ bits.  Qubit and generator locations are not part of this
parameter payload; they are serialized by the complete-token codec below.
\end{definition}

\begin{definition}[Width-aware complete-token codec]
\label{def:token-codec}
Fix a self-delimiting integer code $\mathsf u$, a signed variant
$\mathsf s$, a two-bit Pauli code $\mathsf p$, and a reserved record
terminator $\mathsf d$.  A token $\tau$ with gate-type ID $g$, arity $r$,
ordered operand qubits $q_1,\ldots,q_r$, sparse Pauli support
$\{(z_\ell,\sigma_\ell)\}_{\ell=1}^{s}$, and $k$ parameter slots is encoded
as
\begin{equation*}
\begin{aligned}
\mathsf{ser}_{\mathrm{tok}}(\tau)
={}&\mathsf u(g)\mathsf u(r)
  \prod_{i=1}^{r}\mathsf u(q_i)\;\mathsf u(s)
  \prod_{\ell=1}^{s}
    \bigl(\mathsf u(z_\ell)\mathsf p(\sigma_\ell)\bigr)\\
 &\cdot\mathsf u(k)
  \prod_{a=1}^{k}
    \bigl(\mathsf u(\mathrm{id}_a)
           \mathsf{ser}_{\mathrm{par}}(\lambda_a)\bigr)\mathsf d .
\end{aligned}
\end{equation*}
The zero values of $r$, $s$, or $k$ are explicitly encoded, so the record is
executable and uniquely parseable without out-of-band arity, support, or
delimiter information.  Update, IR-node, and output records use the same
field order, with an opcode specifying which fields are active.

At growing width $n$, an arbitrary qubit or generator choice alone requires
$\lceil\log_2 n\rceil$ bits in any injective codec; the concrete codec above
uses $O(\log n)$ bits per listed address and $O(s\log n)$ bits for an
$s$-sparse Pauli support.  Consequently there is no width-independent
``structural cost per token''.  The literal bit strings
$\mathsf{ser}_{\mathrm{tok}}$ are charged wherever a token is live or
emitted.  Only when $n=n_0$ and a complete-token precision bound $q$ are
both fixed is the token set $\Gamma_{B,n_0}^{\leq q}$ finite; the
fixed-width exact-semantics statements abbreviate it as $\Gamma_B(q)$.
\end{definition}

\begin{definition}[Randomized dynamic-pass compiler and serialized cut budget]
\label{supp:def:compiler}
A compiler $\mathcal A$ may use random-access RAM, stacks, hash tables,
caches, DAG/graph IRs, and intermediate tapes; it may also select a finite
number $p=p(I,\rho)$ of scans or passes dynamically from the input $I$ and
a finite random seed/tape $\rho$.  The external input is append-only on its
first exposure; a new forward pass is re-delivered by the input source from
the start of the stream and is accounted for by the pass count
(\cref{def:w1-pass-accounting}).  Any other access to a consumed
prefix---a copy kept for a later scan through a file, memory mapping,
spool, or cache---is live compiler storage and is charged below.  The final output is write-once.
At every execution cut $c$, serialize a complete restart snapshot in five
disjoint fields:
\begin{itemize}[leftmargin=1.5em]
    \item $\mathsf{ser}_{\mathrm{control}}(c)$: format version, pass and
    finite-control IDs, dynamic pass count and termination/scheduling state,
    all input/output heads and emitted-prefix lengths, end flags, framing
    lengths, and the random seed/tape with its cursor;
    \item $\mathsf{ser}_{\mathrm{store}}(c)$: the persistent and ephemeral
    RAM, stacks, hash-table keys, bucket layout, cache contents and metadata,
    with parameter payloads and IR objects replaced by location tags;
    \item $\mathsf{ser}_{\mathrm{window}}(c)$: the ordered live read-window
    token structures, including the structural fields of
    \cref{def:token-codec} but excluding parameter payloads;
    \item $\mathsf{ser}_{\mathrm{parameter}}(c)$: the location-tagged
    concatenation of every live parameter payload in control, store, window,
    and IR, each encoded by \eqref{eq:precision-charge};
    \item $\mathsf{ser}_{\mathrm{IR}}(c)$: every other live intermediate
    tape/DAG/table object and every rereadable consumed-input block that can
    affect future execution, serialized by node opcode, the structural
    fields of \cref{def:token-codec}, edge endpoints, multiplicity, and
    object boundaries, but with parameter payloads omitted.
\end{itemize}
The corresponding bit-level budget is
\begin{equation}
\label{supp:eq:info-budget}
\boxed{
\begin{aligned}
B_{\mathrm{cut}}(c)
&=B_{\mathrm{control}}(c)+B_{\mathrm{store}}(c)\\
&\quad+B_{\mathrm{window}}(c)+B_{\mathrm{parameter}}(c)
      +B_{\mathrm{IR}}(c).
\end{aligned}}
\end{equation}
where each summand is the literal bit length of the named serialization.
Parameter payloads are the formal angle parameters of live tokens and IR
objects; raw integer data that a compiler computes and keeps, such as
residues or partial sums held in a packed RAM word, belong to
$\mathsf{ser}_{\mathrm{store}}$ and are charged at their literal bit
length.  Every upper bound below uses this convention, and no lower bound
depends on it, because the lower bounds apply to any self-delimiting
serialization whose control field records the emitted-prefix length.
For a fixed run, $B_{\mathrm{cut}}$ is the maximum over its cuts.  For
randomized compilers we use the worst-input expected budget
\begin{equation}
\label{eq:expected-cut}
 \overline B_{\mathrm{cut}}(\mathcal A,\mathcal I)
 :=\max_{I\in\mathcal I}\mathbb E_\rho
       \bigl[\max_c B_{\mathrm{cut}}(c;I,\rho)\bigr].
\end{equation}
No theorem below substitutes a single field, such as
$\max_cB_{\mathrm{parameter}}(c)$ or $\max_cB_{\mathrm{IR}}(c)$, for the
total $B_{\mathrm{cut}}$.  The final output is write-once, is a codeword of
a fixed prefix-free output code, and is charged separately by the output
codes below.
\end{definition}

\begin{remark}[Model scope and escape channels]
\label{rem:escape-channels}
Equation~\eqref{supp:eq:info-budget} charges all restart information, not only
the finite-control state.  In particular, a ZX graph, phase-polynomial
table, spilled update log, hash table, cache, random seed, or materialized
intermediate tape is serialized in the indicated field; formal angle
parameters are serialized in $B_{\mathrm{parameter}}$ and computed integer
data in $B_{\mathrm{store}}$, so every coefficient is charged exactly
once.  Internal rescanning of retained data is allowed,
but the pass schedule, heads, and any retained copy of old input are part
of the snapshot.  A pass re-delivered by the source is accounted by the
pass count rather than by the snapshot, so a bound read off a single cut,
such as \cref{thm:main-tradeoff}, requires the stream to be delivered
once; the transcript bounds of \cref{sec:w1-unified-frontier} cover
re-delivered passes.  The five fields are disjoint by construction, so parameter bits
are not counted again as store, window, or IR structure.  A compiler may
escape a RAM-only workspace lower bound by materializing a large IR, but it
then pays the same bits in $B_{\mathrm{cut}}$.
\end{remark}

\begin{remark}[Physical scope of the bounded-memory streaming model]
\label{rem:streaming-grounding}
The charged streaming model of \cref{supp:def:compiler}---bounded live
state, explicit pass accounting, append-only first exposure---is the
operating regime of a real-time fault-tolerant control stack, not a
convenience abstraction.  Between rounds, the classical electronics
adjacent to the quantum device holds decoder state, Pauli-frame
records, and instruction buffers under hard latency and buffer
budgets: Pauli corrections can be tracked in a classical frame instead of
being applied, which relaxes the timing constraints on syndrome
measurement and decoding~\cite{riesebos2017pauli}, but they must be
resolved before the next non-Clifford gate~\cite{knill2005realistically},
and surface-code compilers are built as streaming pipelines that target
real-time operation~\cite{watkins2024surfacecode}.  Distributed
and modular architectures motivate the same scope from the
interconnect side: a cross-module interaction splits the
\emph{support} of a gate across modules while its coefficient stays
within one module, so modularity bounds the state that crosses a
partition without by itself dispersing any coefficient.  These
architectures therefore ground the model's memory and pass
constraints only; no claim is made that modular partitioning
instantiates the dispersed streams of \cref{sec:physical-dispersed},
whose instances come from scheduling and twirling alone.
\end{remark}

\subsection{Executable Model--Implementation Correspondence}
\label{sec:model-implementation-mapping}

The formal fields above are implemented by one canonical binary codec rather
than by dimension-changing estimates.  The executable specification is
\path{encoding/FORMAT.md}; its reference implementation is
\path{encoding/codec.py}, and the only bit-accounting entry point is
\path{instrumentation/resource_accounting.py}.  A complete codec frame has
the form
\begin{equation}
 \mathtt{UCCACCT0}\,\Vert\,
 \mathsf u(\text{version})\,\Vert\,
 \mathsf u(\text{type})\,\Vert\,
 \mathsf u(|z|)\,\Vert\,z,
 \label{eq:executable-codec-frame}
\end{equation}
where the displayed terminal zero denotes a null byte, $\mathsf u$ is the
canonical base-$128$ unsigned varint, and $z$ is a typed canonical payload.
Signed integers use zigzag coding; rational coefficients are reduced signed
numerator/positive-denominator pairs.  Maps are ordered by their UTF-8 key
bytes.  Qubit, generator, symbol, location, node, edge-endpoint, port, pass,
and head addresses are literal varints in $z$.

\begin{table*}[!t]
\centering
\scriptsize
\setlength{\tabcolsep}{3pt}
\begin{tabularx}{\textwidth}{@{}l l X X@{}}
\toprule
Formal object & Codec object/file & Included information & Measured field \\
\midrule
semantic input & \texttt{SemanticStream} & width and address-complete aggregate records & input artifact bits \\
dispersed flat input & \texttt{FlatUpdateStream} & width, round count, and every additive update & input artifact bits \\
$\mathsf{ser}_{\rm control}$ & \texttt{ControlState}/\texttt{control.uccbin} & pass count/index, heads, scheduling, status, framing and seed state when present & $B_{\rm control}$ \\
$\mathsf{ser}_{\rm store}$ & \texttt{StoreState}/\texttt{store.uccbin} & aggregate tables, maps, caches, stacks and their layout, excluding tagged parameter/IR payloads & $B_{\rm store}$ \\
$\mathsf{ser}_{\rm window}$ & \texttt{LiveWindow}/\texttt{window.uccbin} & ordered live complete-token structure and explicit parameter-slot count & $B_{\rm window}$ \\
$\mathsf{ser}_{\rm parameter}$ & \texttt{ParameterLedger}/\texttt{parameter.uccbin} & sparse symbolic parameters under complete location tuples & $B_{\rm parameter}$ \\
$\mathsf{ser}_{\rm IR}$ & \texttt{IRField}/\texttt{ir.uccbin} & serialized DAG width, opcodes, operands, nodes, edges, ports, boundaries and metadata, with parameters removed & $B_{\rm IR}$ \\
committed prefix & \texttt{committed-prefix.bin} & literal immutable prefix of a codec output frame & $B_{\rm com}$ \\
external output & \texttt{ExternalToolOutput} & tool, format and literal API payload & output artifact bits \\
certificate input & \texttt{CertificateInput} & certificate kind, certified-artifact SHA-256 and literal sidecar & certificate artifact bits \\
pass count & manifest \texttt{passes.count} and row \texttt{p} & actual declared scans of the instrumented compiler & $p$ \\
\bottomrule
\end{tabularx}
\caption{One-to-one map from formal serialized objects to executable artifacts.
No row permits a gate, token, support, or IR-node count to stand in for bits.}
\label{tab:model-implementation-map}
\end{table*}

\begin{lemma}[Canonical prefix-free implementation]
\label{lem:codec-prefix-free}
Version-one complete frames are injective and prefix-free.  Concatenations of
frames are uniquely parseable without an out-of-band arity or record count.
The semantic and flat decoders preserve their exact aggregate coefficient
table.
\end{lemma}

\begin{proof}
Canonical varints have a unique terminating byte, and the version, type, and
payload length are therefore uniquely decoded.  The declared payload length
locates the end of the frame.  A complete valid frame cannot be a proper prefix
of another complete valid frame: the shorter declared payload ends at its frame
boundary, whereas any following byte belongs to a subsequent frame.  Typed
payload values have distinct one-byte tags; lists and maps carry counts, maps
have strictly increasing keys, and reduced rational pairs are unique.  This
proves injectivity and unique concatenated parsing.  Finally,
\texttt{semantic\_signature} sums exact rational sparse parameters by the
literal pair (generator address, Pauli support), so encode--decode leaves the
aggregate table unchanged for both input representations.
\end{proof}

\begin{proposition}[File-backed accounting soundness]
\label{prop:file-backed-accounting}
Suppose the version-one verifier accepts a cut manifest $M$.  For every
$f\in\{\mathrm{control},\mathrm{store},\mathrm{window},
\mathrm{parameter},\mathrm{IR},\mathrm{com}\}$, let $A_f$ be the file named by
$M$.  Then the implementation reports
\begin{equation}
\begin{aligned}
 B_f&=8|A_f|_{\rm byte},\\
 B_{\rm cut}&=B_{\rm control}+B_{\rm store}+B_{\rm window}\\
             &\quad+B_{\rm parameter}+B_{\rm IR},
\end{aligned}
 \label{eq:file-backed-accounting}
\end{equation}
and reports the manifest pass count as $p$.  The reported values change when a
serialized address crosses a varint-width boundary.  They cannot be obtained
from a node or gate count alone.
\end{proposition}

\begin{proof}
The accountant first writes the five complete frames and the literal committed
prefix, then obtains every byte count by \texttt{stat}, stores the SHA-256 of
each file, writes $M$, and invokes the verifier before returning a row.  The
verifier reopens all six files, recomputes sizes and digests, decodes the five
restart fields, and rejects any disagreement.  It then forms the sum in
\eqref{eq:file-backed-accounting}; there is no count-to-bit code path.  Since an
address is a literal canonical varint, for example $127$ and $128$ occupy
different address lengths.  Conversely, two one-node DAGs may contain those
different addresses and hence have unequal file lengths despite equal node
count.
\end{proof}

\begin{remark}[Observable boundary versus restart cut]
\label{rem:observable-boundary-accounting}
The masked-share tracer serializes the five declared cut fields after every
update.  The verifier reopens and decodes every recorded field slice, checks
its size and digest, and recomputes the cut sum; its maximum is therefore a
file-backed measurement of the declared $B_{\rm cut}$ coordinate.  The W3
campaign does not claim that the Python process is destroyed and restarted
between events.  External compilers do not expose their private in-process DAGs,
caches, or tapes.  Their matched-matrix and natural-factorial rows consequently
report a file-backed \emph{API-boundary} snapshot and are diagnostic only; they
are not promoted to measurements of the maximum private $B_{\rm cut}$.  Peak
RSS remains a separate byte dimension and is not substituted for any field in
\eqref{supp:eq:info-budget}.
\end{remark}

\begin{remark}[Self-contained and conditional descriptions]
\label{rem:dictionary-accounting-code}
Both input-stream objects declare either \texttt{self\_contained} or
\texttt{conditional}.  The former forbids an external dictionary.  The latter
requires the SHA-256 of the exact public family dictionary, and its manifest
records the dictionary's separate file size and whether it is charged at the
cut.  Thus a conditional description cannot silently be reported as a
self-contained circuit code.
\end{remark}


\subsection{Restart from the serialized cut}

\begin{lemma}[Restart from the serialized cut]
\label{app:lem:reconstruction}
Fix $\mathcal A$ as in \cref{supp:def:compiler}, receiving a stream of
\cref{def:w1-dispersed-stream} once.  At the crossing $c_{A\to B}$ before
the last round, the five fields in \eqref{supp:eq:info-budget}, together
with any prefix already committed to the final write-once output, determine
the continuation of that run for every supplied last round $v$.  For a randomized run this statement is conditional on the serialized
seed/tape and cursor.
\end{lemma}

\begin{proof}
The control field restores the exact transition, all head positions,
dynamic pass/termination schedule, and random seed/tape cursor.  The store
and window fields restore RAM, stack, hash/cache metadata, nonparameter
machine data, and ordered live token structures.  The parameter field
restores every omitted coefficient by its location tag.  The IR field
restores every other live intermediate tape, DAG, table, addressed record,
and any consumed input block retained for rescanning.  These disjoint fields
therefore reconstruct the complete machine configuration.  Conditional on
the serialized random tape, the receiver can append the supplied second
round and execute the same dynamically chosen continuation.  The separately
transmitted final-output prefix is prepended to the suffix produced after
restart.  No additional state, old-input file, or public-family side channel
is used.  The single-delivery hypothesis is needed: if the source
re-delivered the stream for a later pass, the continuation would reread the
first round, which the receiver does not hold.
\end{proof}
